\section{Design Decisions and Issues} \label{sec:design-decisions-and-issues}

\subsection{General CNF Instead of 3-CNF} \label{subsec:general-cnf-instead-of-3-cnf}

The encoding uses general CNF clauses rather than forcing every clause to contain exactly three literals.
Although 3-SAT is important theoretically, converting every constraint into 3-CNF would require additional auxiliary variables and would make the implementation less readable.
General CNF is accepted directly by the PySAT solver interface, so this choice keeps the encoding simpler.

\subsection{Pairwise Cardinality Encoding} \label{subsec:pairwise-cardinality-encoding}

The system uses pairwise clauses for at-most-one constraints.
This produces $O(r^2)$ clauses for a group of $r$ literals, but the implementation is straightforward and works well for the tested problem sizes.

\subsection{Iterative Deepening for Optimality} \label{subsec:iterative-deepening-for-optimality}

Finding the shortest plan is achieved by repeatedly solving fixed-horizon instances.
This approach is easy to implement and gives a strong optimality guarantee: the first satisfiable horizon is the shortest possible one.

\subsection{Readable Mapping Output} \label{subsec:readable-mapping-output}

The project explicitly exports a mapping from integer SAT variables to readable names.
This is useful both for debugging and for satisfying the requirement that the SAT instance and truth assignment be interpretable.

\subsection{Frontend Visualization} \label{subsec:frontend-visualization}

The frontend is not required for the logical encoding itself, but it makes the solver much easier to understand.
Instead of only returning a list of Boolean assignments, the system shows a step-by-step animation of the plan.